%% notes-tex/025-graph-reg-sweep/sections/1-question.tex

\section{The question\secstatus{todo}}
\label{sec:question}

The cell graph transformer of experiment 010 attends over all 6{,}607 genes in every
layer. Nine gene-gene graphs, physical and regulatory interactions, TFLink, and the six
STRING channels (neighborhood, fusion, co-occurrence, co-expression, experimental,
database), are given to nine heads of layer 1, one graph per head. Two mechanisms can
hand a graph to a head. The \term{soft prior} adds to the training loss, for each head
$k$ with row-normalized adjacency $\tilde{A}_k$ and attention $\alpha_k$ on the
wild-type graph,
\begin{equation}
  \lambda \sum_{k=1}^{9} \frac{N}{|E_k|}
  \sum_{i} D_{\mathrm{KL}}\!\left(\tilde{A}_{k,i\cdot} \,\|\, \alpha_{k,i\cdot}\right),
  \label{eq:kl}
\end{equation}
with $N$ the number of genes and $|E_k|$ the edge count of graph $k$, so a head is pulled
toward spreading each gene's attention uniformly over that gene's neighbors and is free
to put mass anywhere. The \term{hard mask} sets the attention logits of every
non-edge (self-loops kept) to $-\infty$ before the softmax in layer 1, so a head can only
attend along the graph and needs no penalty and no $\lambda$. The two are the ends of one
axis: the mask is the $\lambda \to \infty$ limit of Eq.~\ref{eq:kl} on the support, but
not in the weights, since the mask lets a head weight its neighbors as it likes while the
prior pulls toward the uniform row of $\tilde{A}_k$.

Earlier readings put the two mechanisms far apart. The 010 model at $\lambda = 10^{-3}$
under a cosine schedule reached 0.456 to 0.464 held-out Pearson over three seeds, and on
the 025 build the same configuration replicated at 0.446 (GilaHyper job 1598, W\&B
\texttt{0yw7moue}), while the two hard-mask arms run then collapsed: 0.307 at epoch 1
and then near 0 (\texttt{7f1yrsq9}), or near 0 for 17 epochs (\texttt{4qmgkcgn}). Those
mask runs were under the cosine schedule with a different normalizer, so the collapse was
not a controlled comparison. The hypothesis this sweep was designed to test is stated in
the form it was proposed: \emph{the graphs help, but only when the model is optimized
through the divergence; imposing them as a mask does not help}. A second question rides
on the first: if the prior helps, is it the biology of the nine graphs or would any graph
with the same degree sequence condition training the same way?

\notesec{Design}
Every arm is the constant-rate control \texttt{cgt\_s0\_r\_kl\_ctrl\_013} with one key
changed (Table~\ref{tab:design}). The control is the 376{,}732 triples of experiment 010
(subset S0 of the 025 build) on 010's pinned random split, 301{,}386 training and
37{,}673 validation triples, a learnable 1{,}000-dimensional gene table, the label
normalizer fit on the training split, AdamW at $2.5 \times 10^{-4}$ with no schedule,
batch 256 over four GPUs, 30 epochs. The runs are Delta \texttt{gpuA40x4} jobs, one node
and four A40s per run, 24 h clocks, the 554 GB LMDB staged onto node-local NVMe at the
start of each job (\file{experiments/025-solid-growth/scripts/delta_cgt.slurm},
\file{experiments/025-solid-growth/scripts/delta_submit_sweep.sh}; jobs 22055147 to
22055169, 22056267 to 22056269, 22080055, 22080057 and 22107791). The prior arms peak at
6.05 GB of allocated GPU memory per card, the mask at 2.76 GB, because only the
penalized layer materializes its attention.

%% SOURCE: hand-authored from experiments/025-solid-growth/scripts/delta_submit_sweep.sh and the
%% configs cgt_s0_r_kl_ctrl_013, cgt_s0_r_mask_028, cgt_s0_r_kl_rand_031; seed counts from
%% experiments/025-solid-growth/results/graph_reg_sweep_summary.json (graph_reg_sweep_readout.py).
\begin{table}[htbp]
\centering
\footnotesize
\caption{The arms. Every arm is \texttt{cgt\_s0\_r\_kl\_ctrl\_013} with the one change
named; the control itself is the ladder's $\lambda = 10^{-3}$ point. Complete runs are
those with 30 logged epochs on W\&B at the pull of 2026-09-27; a seed listed as partial
was still running.}
\label{tab:design}
\begin{tabular}{llll}
\toprule
arm & change on the control & graph target & complete seeds \\
\midrule
no penalty & \texttt{graph\_reg\_lambda=0} & none & 1, 2, 3 \\
prior ladder & \texttt{graph\_reg\_lambda} $\in \{10^{-5}, 10^{-4}, 10^{-2}, 10^{-1}, 1\}$ & nine biological graphs, layer 1 & 1, 2, 3 each; $10^{-5}$: 1, 3 \\
control & none & nine biological graphs, layer 1, $\lambda = 10^{-3}$ & 1, 2, 3 \\
hard mask & \texttt{cgt\_s0\_r\_mask\_028}: mask on, $\lambda = 0$ & nine biological graphs, layer 1 & 1, 2, 3 \\
random graphs & \texttt{cgt\_s0\_r\_kl\_rand\_031}: rewiring on & nine rewired graphs, layer 1, $\lambda = 10^{-3}$ & 1, 2; 3 partial (epoch 16) \\
\bottomrule
\end{tabular}
\end{table}

\notesec{Readings}
Every run is read twice. The \term{fixed-epoch reading} is validation Pearson at epoch
29, the protocol's last epoch, the same optimization state for every arm. The
\term{max-over-epochs reading} is the largest validation Pearson over the 30 epochs, an
upward-biased order statistic reported with the epoch at which it occurred; it is the
number a best-checkpoint deployment would see. Arm values are the mean and sample
standard deviation over seeds; a paired comparison is by seed, since the split is pinned
and the seed sets the initialization and the batch order for every arm alike.

\notesec{Vocabulary}
Terms used in the figures and tables, in alphabetical order.

%% SOURCE: hand-authored; definitions follow torchcell/models/equivariant_cell_graph_transformer.py
%% (compute_graph_regularization_loss) and torchcell/trainers/int_transformer_cell.py (gradient probe,
%% edge recovery).
\begin{table}[htbp]
\centering
\footnotesize
\caption{Controlled vocabulary.}
\label{tab:vocab}
\begin{tabular}{p{3.2cm}p{13.6cm}}
\toprule
term & meaning \\
\midrule
divergence & Eq.~\ref{eq:kl} without its $\lambda$: the logged penalty divided by $\lambda$, summed over the nine heads and normalized by edge count, computed on the wild-type graph. Comparable across the ladder; not logged at $\lambda = 0$ or under the mask, where the model skips the computation. Reported from the validation pass at epoch 29. \\
edge recall at degree & For each gene $i$ of graph $k$ with degree $d_i$, the fraction of $i$'s neighbors among the $d_i$ largest entries of $\alpha_{k,i\cdot}$; averaged over genes, then over the nine heads. It is 1 by construction under the mask and near $1/N$ for attention unrelated to the graph. Logged at epochs 9, 19 and 29. \\
gradient probe & On the first batch of epochs 0, 1, 2, 5, 10 and 20, one extra backward pass per loss term; the $\ell_2$ norm of the parameter gradient of the point loss and of the penalty, before clipping, and their ratio. \\
point loss & Mean squared error of the prediction against the z-scored interaction label. \\
random graphs, degree-matched & Each of the nine graphs rewired by seeded double-edge swaps (five swaps per edge) that keep every gene's degree, in and out for the directed graphs; the rewiring seed is the run seed. The prior and the edge diagnostics are computed against the rewired graphs. \\
S0 & The 376{,}732 measured triples of experiment 010, as stored in the 025 build; the pool of every arm here. \\
\bottomrule
\end{tabular}
\end{table}
